\section*{第 \ref{ch:b3:lebesgue} 章 --- 勒贝格
积分}

\begin{solution}{exo:b3:lebesgue:1}
(a) 若 $f$ 非减，则 $\{f > t\}$ 为 $\varnothing$、$\R$，
或为射线 $\intoo a{+\infty}$ / $\intco a{+\infty}$，在
各情形下均为博雷尔集；递减情形同理。导数 $f'(x) =
\lim_n n\bigl(f(x + \frac1n) - f(x)\bigr)$ 是
连续（从而可测）函数的逐点极限：
由\cref{prop:b3:lebesgue:stability}(c)可知。

(b) $\{f > t\} = \bigcup_{q \in \Q,\, q > t}\{f > q\}$：若
有理水平集可测，则一切水平集皆可测，且射线
生成 $\mathcal B(\R)$。
\end{solution}

\begin{solution}{exo:b3:lebesgue:2}
其一：$(1 + x/n)^n \nearrow \eu^x$（$x \geq 0$），故
被积函数逐点趋于 $\eu^{-x}\cos x$；当 $n \geq 2$ 时，
$(1 + x/n)^n \geq (1 + x/2)^2$，由此得可积控制函数
$(1 + x/2)^{-2}$。由控制收敛定理：
\[
\lim_n\int_0^\infty\frac{\cos x}{(1 + x/n)^n}\dd x
= \int_0^\infty \eu^{-x}\cos x\,\dd x
= \operatorname{Re}\int_0^\infty\eu^{-(1 - \iu)x}\dd x
= \operatorname{Re}\frac{1}{1 - \iu} = \frac12 .
\]
其二：作代换 $u = x^n$（为 $\mathcal C^1$
$\intoo01$ 上的双射）：
\[
\int_0^1\frac{nx^{n-1}}{1 + x}\dd x = \int_0^1\frac{\dd u}{1 +
u^{1/n}} \longrightarrow \int_0^1\frac{\dd u}{2} = \frac12,
\]
由控制收敛定理：对 $u \in \intoo01$，$u^{1/n} \to 1$，且
被积函数在有限测度空间上以 $1$ 为界。
\end{solution}

\begin{solution}{exo:b3:lebesgue:3}
(a) 取 $f_n = n\,\mathbf 1_{\intoo0{1/n}}$：逐点有 $\liminf f_n = 0$，
$\int f_n = 1$，故 $0 < 1$。

(b) 高度：$n\mathbf 1_{\intoo0{1/n}}$；宽度：$\frac1n\mathbf
1_{\intoo0n}$；平移：$\mathbf 1_{\intoo n{n+1}}$。全部
逐点趋于 $0$，且 $\int = 1$。各情形下
\emph{控制}假设不成立：$\sup_nf_n$ 为
$\approx 1/x$（在 $0$ 附近）、$\approx$ 不可积的
常值型，形如 $\mathbf 1_{\intoo1\infty}$ 者 --- 皆不可积。

(c) “$\int\limsup f_n \geq \limsup\int f_n$” 不成立，
对平移凸峰：左边为 $0$，右边为 $1$。
“$\limsup\int \leq \int\limsup$”是同一结论。又
对 $\limsup$ 的法图引理，把 $\leq$ 反向
（“逆法图引理”）需要一个控制函数 --- 同一凸峰
即为反例。
\end{solution}

\begin{solution}{exo:b3:lebesgue:4}
(a) 当 $x > 0$ 时：$\frac1{\eu^x - 1} = \frac{\eu^{-x}}{1 -
\eu^{-x}} = \sum_{n\geq1}\eu^{-nx}$，故 $\frac{x}{\eu^x - 1} =
\sum_{n\geq1}x\eu^{-nx}$，为非负可测函数的级数，
由\cref{cor:b3:lebesgue:additivity}，可逐项
积分：
\[
\int_0^\infty\frac{x\,\dd x}{\eu^x - 1}
= \sum_{n\geq1}\int_0^\infty x\eu^{-nx}\dd x
= \sum_{n\geq1}\frac1{n^2} = \frac{\pi^2}6
\]
（$\int_0^\infty x\eu^{-nx}\dd x = n^{-2}$，分部积分可得；巴塞尔
问题见本科二年级分册，或见后文\cref{exo:b3:hilbert:5}）。

(b) 在 $\intoo01$ 上，$-x\ln x \geq 0$，故 $x^{-x} =
\eu^{-x\ln x} = \sum_k\frac{(-x\ln x)^k}{k!}$ 为
非负项的级数，再次交换求和与积分。作代换 $x =
\eu^{-u/(k+1)}$：
\[
\int_0^1(-x\ln x)^k\dd x
= \int_0^\infty\Bigl(\frac{u}{k+1}\Bigr)^{k}
\eu^{-\frac{ku}{k+1}}\;\frac{\eu^{-\frac u{k+1}}}{k+1}\,\dd u
= \frac{1}{(k+1)^{k+1}}\int_0^\infty u^k\eu^{-u}\dd u
= \frac{k!}{(k+1)^{k+1}} .
\]
于是 $\int_0^1x^{-x}\dd x = \sum_{k\geq0}\frac{1}{(k+1)^{k+1}}
= \sum_{n\geq1}n^{-n}$：即“大二学生之梦”，严格成立。
\end{solution}

\begin{solution}{exo:b3:lebesgue:5}
(a) 马尔可夫不等式：$a\,\mathbf 1_{\{f \geq a\}} \leq f$，积分得
$\mu(\{f \geq a\}) \leq \frac1a\int f$。若 $\int f = 0$：
$\mu(\{f \geq \frac1n\}) = 0$（对每个 $n$），且 $\{f > 0\} =
\bigcup_n\{f \geq \frac1n\}$ 为零测集。

(b) $\mu(\{\abs f = \infty\}) \leq \mu(\{\abs f \geq n\}) \leq
\frac1n\int\abs f \to 0$.

(c) 取 $A = \{f > 0\}$：$\int f^+\dd\mu = \int_Af\,\dd\mu =
0$，故 $f^+ = 0$ 几乎处处（由 (a)）；同理 $f^- = 0$ 几乎处处。
\end{solution}

\begin{solution}{exo:b3:lebesgue:6}
(a) $\mathbf 1_\Q$：处处不连续，
不是黎曼可积的（见\cref{thm:b3:lebesgue:riemann}）；其
勒贝格积分为 $\lambda(\Q) = 0$。托马函数在
每个无理点处连续（给定 $\varepsilon$，
$\intcc01$ 中分母 $\leq
1/\varepsilon$ 的有理数只有有限个；取小邻域避开它们），
在有理点处不连续（无理数稠密）：
几乎处处连续，黎曼可积，积分为 $0$（它几乎处处为零）。
$\mathbf 1_K$，$K$ 为胖康托尔集：间断点集为
$\partial K = K$（闭集且内部为空），测度为
$\frac12 > 0$：\emph{不}黎曼可积 --- 然而
勒贝格可积，积分为 $\lambda(K) = \frac12$。

(b) $\int_{k\pi}^{(k+1)\pi}\frac{\abs{\sin x}}x\dd x \geq
\frac1{(k+1)\pi}\int_{k\pi}^{(k+1)\pi}\abs{\sin x}\dd x =
\frac{2}{(k+1)\pi}$，故级数发散。反常积分
的收敛性：对 $A > \pi$，
\[
\int_\pi^A\frac{\sin x}x\dd x = \Bigl[\frac{-\cos
x}x\Bigr]_\pi^A - \int_\pi^A\frac{\cos x}{x^2}\dd x,
\]
当 $A \to \infty$ 时两项均收敛（$\frac1{x^2}$ 可积）：
即条件收敛，
而非绝对可积。
\end{solution}

\begin{solution}{exo:b3:lebesgue:7}
控制：$\abs{\partial_t(\eu^{-x^2}\cos(tx))} =
\abs{x\eu^{-x^2}\sin(tx)} \leq x\eu^{-x^2}$，可积且
与 $t$ 无关，故由\cref{thm:b3:lebesgue:paramdiff}
可整体应用，
\[
F'(t) = -\int_0^\infty x\eu^{-x^2}\sin(tx)\,\dd x
= \Bigl[\tfrac12\eu^{-x^2}\sin(tx)\Bigr]_0^\infty
- \frac t2\int_0^\infty\eu^{-x^2}\cos(tx)\dd x
= -\frac t2F(t)
\]
（取 $\dd v = x\eu^{-x^2}\dd x$，作分部积分）。
线性常微分方程给出 $F(t) = F(0)\eu^{-t^2/4}$，且 $F(0) =
\frac{\sqrt\pi}2$（见\cref{pb:b3:lebesgue:1}），高斯函数
在此余弦变换下自我再现。
\end{solution}

\begin{solution}{exo:b3:lebesgue:8}
积分收敛：在 $0$ 附近被积函数趋于 $b - a$
（有界），且在无穷远处以 $\eu^{-ax}$ 的速度衰减。取定 $a$
并记 $I(b) = \int_0^\infty\frac{\eu^{-ax} -
\eu^{-bx}}x\dd x$，将其视为 $b \in \intco a{+\infty}$ 上的函数。
对所有 $b \geq a$：$\abs{\partial_b(\text{被积函数})} =
\eu^{-bx} \leq \eu^{-ax}$，且 $\int_0^\infty\eu^{-ax}\dd x =
\frac1a < \infty$：可积控制函数存在。于是
由\cref{thm:b3:lebesgue:paramdiff}可知 $I'(b) =
\int_0^\infty\eu^{-bx}\dd x = \frac1b$，且 $I(a) = 0$：
\[
I(b) = \int_a^b\frac{\dd t}t = \ln\frac ba .
\]
（等价地，按提示的路线，被积函数为
$\int_a^b\eu^{-xt}\dd t \geq 0$，而交换次序恰是
\cref{cor:b3:lebesgue:additivity}的连续类比，
即托内利定理 --- 在\cref{ch:b3:product}中证明；
而参数化的路线则完全在本章之内。）
\end{solution}

\begin{solution}{exo:b3:lebesgue:9}
$\nu(\varnothing) = 0$；对两两不交的 $(A_n)$，$f\mathbf
1_{\bigsqcup A_n} = \sum_nf\mathbf 1_{A_n}$（逐点成立，各项
$\geq 0$），由\cref{cor:b3:lebesgue:additivity}可得
$\sigma$-可加性。公式 $\int g\,\dd\nu = \int
gf\,\dd\mu$：对 $g = \mathbf 1_A$，即为
$\nu$ 的定义；对简单函数 $g$，由线性性；对可测的 $g \geq 0$，
取简单函数 $s_n \nearrow g$
（见\cref{thm:b3:lebesgue:approximation}）：$s_nf \nearrow gf$，且
两侧由单调收敛定理取极限。（此“示性函数
$\to$ 简单函数 $\to$ 单调收敛定理”这一递进路线是
理论中的标准机制。）
\end{solution}

\begin{solution}{exo:b3:lebesgue:10}
(a) $\abs{\sin(tx)} \leq \abs tx$ 给出
$\abs{\frac{\sin(tx)}{x(1+x^2)}} \leq \frac{\abs
t}{1+x^2}$：积分收敛，且在 $\abs t \leq T$ 上
控制函数 $\frac{T}{1+x^2}$ 给出连续性
（见\cref{thm:b3:lebesgue:paramcont}）。求导：
$\abs{\partial_t} = \abs{\frac{\cos(tx)}{1+x^2}} \leq
\frac1{1+x^2}$，可积：$f'(t) =
\int_0^\infty\frac{\cos(tx)}{1+x^2}\dd x$ 对所有 $t$ 成立。若再求一次导，
则需对 $\frac{x\sin(tx)}{1 +
x^2}$ 积分，其绝对值行为类似 $\frac{\abs{\sin(tx)}}
x$（在无穷远处）：\emph{不}可积 --- 不存在控制函数，
因而 \cref{thm:b3:lebesgue:paramdiff} 不能再次应用。

(b) 承认 $f'(t) = \frac\pi2\eu^{-t}$ 对 $t > 0$ 成立：由
$f$ 为奇函数，$f'$ 为偶函数，故 $f'(t) =
\frac\pi2\eu^{-\abs t}$ --- 它在 $0$ 处\emph{不可导
}：$f$ 属于 $\mathcal C^1$ 但不属于 $\mathcal C^2$，
这说明被阻断的二阶求导并非
技术性的偶然。对 $t > 0$，反常积分
$\int_0^\infty\frac{x\sin(tx)}{1+x^2}\dd x$ 等于 $-f''(t) =
\frac\pi2\eu^{-t}$（对承认的公式在合法之处求导，
即在 $\intoo0\infty$ 上） --- 一个反常的，
非勒贝格意义下的值。
\end{solution}

\begin{solution}{exo:b3:lebesgue:11}
(a) 记 $g_n = (f - f_n)^+$，则 $0 \leq g_n \leq f$
（由 $f_n$ 非负），$g_n \to 0$ 几乎处处成立，且 $f$ 为
可积控制函数：$\int g_n \to 0$（由控制收敛定理）。注意到
$\abs{f_n - f} = 2(f - f_n)^+ - (f - f_n)$,
\[
\int\abs{f_n - f} = 2\int g_n - \Bigl(\int f - \int
f_n\Bigr) \longrightarrow 0 + 0 .
\]

(b) 滑动小峰 $f_n = \mathbf 1_{\intcc n{n+1}}$ 具有
$f_n \to 0$ 几乎处处，且 $\int f_n = 1 \not\to 0$：若无
积分的收敛性，$L^1$ 收敛失效（
假设亦失效）。带符号的例子：$f_n \to 0$ 在
每个 $x \neq 0$ 处，$\int f_n = 0$，但 $\int\abs{f_n} = 2$：
对带符号的序列，该定理实质上关于
$\abs{f_n}$ 型的控制，而非负性恰好用在
$g_n \leq f$ 处。

(c) 密度满足 $\int p_n = 1 = \int p$；
(a) 的假设自动成立，故 $p_n \to p$ 几乎处处蕴含
$\norm{p_n - p}_{L^1} \to 0$ --- 从而概率
$\int_Ap_n \to \int_Ap$ 对所有
可测的 $A$ \emph{一致}成立（$\abs{\int_A(p_n - p)} \leq
\norm{p_n - p}_1$）：谢费把密度的逐点收敛
转化为分布的全变差收敛。
\end{solution}

\begin{solution}{exo:b3:lebesgue:12}
(a) 在 $\intoo0n$ 上，$\varphi_n(x) = n\log(1 - \frac xn)$
关于 $n$ 递增并趋于 $-x$（映射 $t \mapsto
\frac{\log(1 - xt)}{t}$ 随 $t = \frac1n
\downarrow 0$ 时递减；或直接展开，$\varphi_{n}' \geq 0$ 对 $n$ 成立，借助
$\log(1-u) + \frac{u}{1-u} \geq 0$）。于是 $(1 -
\frac xn)^n\eu^{x/2}\mathbf 1_{x<n} \nearrow
\eu^{-x/2}$，由单调收敛定理得
\[
\lim_n\int_0^n\Bigl(1 - \frac xn\Bigr)^n\eu^{x/2}\dd x =
\int_0^\infty\eu^{-x/2}\dd x = 2 .
\]

(b) 作代换 $u = x^n$（$x = u^{1/n}$，$n x^{n-1}\dd x =
\dd u$）：
\[
\int_0^1\frac{nx^{n-1}}{1 + x}\dd x =
\int_0^1\frac{\dd u}{1 + u^{1/n}} .
\]
当 $u \in \intoo01$ 时，$u^{1/n} \to 1$，故被积函数
趋于 $\frac12$，被 $1$ 控制，极限为
$\frac12$（由控制收敛定理）。（$nx^{n-1}$ 的质量集中在 $x
= 1$ 处，即 $\frac1{1+x} = \frac12$：换元使
集中现象清晰可见。）

(c) 逐点地，$(1 + x/n)^n \nearrow \eu^x$ 且 $x^{1/n} \to
1$（$x > 0$），被积函数趋于 $\eu^{-x}$。控制函数
当 $n \geq 2$ 时：在 $\intoc01$ 上，$x^{-1/n} \leq x^{-1/2}$ 且
$(1 + x/n)^{-n} \leq 1$：控制为 $x^{-1/2}$，可积；在
$\intoo1\infty$ 上，$x^{-1/n} \leq 1$ 且 $(1 + x/n)^n \geq 1 +
\binom n2\frac{x^2}{n^2} \geq 1 + \frac{x^2}4$：控制为
$\frac{4}{4 + x^2}$，可积。由控制收敛定理得：
\[
\lim_n\int_0^\infty\frac{\dd x}{(1 + x/n)^nx^{1/n}} =
\int_0^\infty\eu^{-x}\dd x = 1 .
\]
\end{solution}

\begin{solution}{pb:b3:lebesgue:1}
\textbf{1.} $A$ 是 $\mathcal C^1$ 的，这由微积分基本定理及
链式法则可得：$A'(t) =
2\eu^{-t^2}\int_0^t\eu^{-x^2}\dd x$。对 $B$：
$\partial_t\bigl(\frac{\eu^{-t^2(1+x^2)}}{1+x^2}\bigr) =
-2t\,\eu^{-t^2(1+x^2)}$，在 $[t_0, T]
\times \intcc01$ 上连续且有界，对任意 $0 < t_0 < T$ 成立（有限
测度空间上的有界控制即可）：$B$ 是 $\mathcal C^1$ 函数，
在 $\intoo0{+\infty}$ 上，其导数为
\[
B'(t) = -2t\int_0^1 \eu^{-t^2(1+x^2)}\dd x
= -2\eu^{-t^2}\int_0^1 t\,\eu^{-(tx)^2}\dd x
= -2\eu^{-t^2}\int_0^t\eu^{-u^2}\dd u = -A'(t)
\]
（作代换 $u = tx$）。

\textbf{2.} $A(0) = 0$，且 $B(0) = \int_0^1\frac{\dd x}{1+x^2}
= \frac\pi4$。$A + B$ 的导数在
$\intoo0{+\infty}$ 上为零，且在 $0$ 处连续（$B$ 由
以 $\frac1{1+x^2}$ 为控制，并
借助\cref{thm:b3:lebesgue:paramcont}）：$A + B \equiv \frac\pi4$。
当 $t\to+\infty$ 时：$0 \leq B(t) \leq \eu^{-t^2} \to 0$，且
$A(t) \to \bigl(\int_0^\infty\eu^{-x^2}\dd x\bigr)^2$（由单调收敛定理，或
仅对内层积分用单调收敛）。

\textbf{3.} 因此 $\bigl(\int_0^\infty\eu^{-x^2}\dd x\bigr)^2
= \frac\pi4$，故 $\int_0^\infty\eu^{-x^2}\dd x =
\frac{\sqrt\pi}2$，由偶性得 $G = \sqrt\pi$。又
$\Gamma(\frac12) = \int_0^\infty x^{-1/2}\eu^{-x}\dd x
\overset{x = u^2}{=} 2\int_0^\infty\eu^{-u^2}\dd u =
\sqrt\pi$.

\textbf{4.} 当 $t > 0$ 时：$\abs{\eu^{-tx}\frac{\sin x}x} \leq
\eu^{-tx}$，可积。当 $t = 0$ 时需考察广义收敛：在
$[\pi, A]$ 上，
\[
\int_\pi^A\frac{\sin x}x\dd x
= \Bigl[-\frac{\cos x}x\Bigr]_\pi^A -
\int_\pi^A\frac{\cos x}{x^2}\dd x,
\]
两项当 $A \to \infty$ 时均收敛；在 $0$ 附近，被积函数
可连续延拓为 $1$。

\textbf{5.} 在 $[t_0, +\infty)$ 上（$t_0 > 0$）：
$\abs{\partial_t\bigl(\eu^{-tx}\tfrac{\sin x}x\bigr)} =
\eu^{-tx}\abs{\sin x} \leq \eu^{-t_0x}$，可积：
\cref{thm:b3:lebesgue:paramdiff} 在每个此类
区间上均适用，因此在整个 $\intoo0{+\infty}$ 上：
\[
F'(t) = -\int_0^\infty\eu^{-tx}\sin x\,\dd x
= -\operatorname{Im}\int_0^\infty\eu^{-(t - \iu)x}\dd x
= -\operatorname{Im}\frac{1}{t - \iu} = -\frac{1}{1 + t^2}.
\]

\textbf{6.} $\abs{F(t)} \leq \int_0^\infty\eu^{-tx}\dd x =
\frac1t \to 0$。积分 $F' = -\frac1{1+t^2}$ 得 $F(t) = C -
\arctan t$，由 $t \to \infty$ 知 $C = \frac\pi2$：$F(t) =
\frac\pi2 - \arctan t$，在 $\intoo0{+\infty}$ 上成立。

\textbf{7.} 在 $[A, R]$ 上分部积分，取 $\sin x =
(-\cos x)'$，并令 $R \to \infty$：
\[
\int_A^{\infty}\eu^{-tx}\frac{\sin x}x\dd x
= \frac{\cos A\;\eu^{-tA}}{A}
- \int_A^\infty \cos x\,\Bigl(\frac tx +
\frac1{x^2}\Bigr)\eu^{-tx}\dd x .
\]
放缩 $\abs{\cos} \leq 1$：第一项 $\leq \frac1A$；
积分至多为 $\int_A^\infty t\eu^{-tx}\frac{\dd x}x
+ \int_A^\infty\frac{\dd x}{x^2} \leq \frac1A\int_0^\infty
t\eu^{-tx}\dd x + \frac1A = \frac2A$。合计 $\leq \frac 3A$，
对 $t \in [0, 1]$ 一致成立（包括 $t = 0$ 的情形）。现在
\[
\abs{F(t) - D} \leq
\Bigl|\int_0^A(\eu^{-tx} - 1)\,\frac{\sin x}x\,\dd x\Bigr| +
\frac6A .
\]
在 $[0, A]$ 上：$\abs{\eu^{-tx} - 1} \leq tx \leq tA$，且
$\abs{\frac{\sin x}x} \leq 1$，故第一项至多为
$tA^2$。取 $A$ 使 $\frac6A < \varepsilon$，再取 $t <
\varepsilon/A^2$，则 $\abs{F(t) - D} < 2\varepsilon$。

\textbf{8.} 因此 $D = \lim_{t\to0^+}F(t) =
\lim_{t\to0^+}\bigl(\frac\pi2 - \arctan t\bigr) =
\frac\pi2$。

\textbf{9.} 分部积分（$u = \sin^2x$，$v' = x^{-2}$）：
\[
\int_0^\infty\frac{\sin^2x}{x^2}\dd x
= \Bigl[-\frac{\sin^2x}{x}\Bigr]_0^\infty +
\int_0^\infty\frac{2\sin x\cos x}{x}\dd x
= \int_0^\infty\frac{\sin 2x}{x}\dd x = D = \frac\pi2
\]
（末步作代换 $u = 2x$；边界项为零：
$\sin^2 x/x \to 0$，两端皆然）。

\textbf{10.} 分部积分（$u = 1 - \cos x$，$v' = x^{-2}$）：
$\int_0^\infty\frac{1 - \cos x}{x^2}\dd x =
\int_0^\infty\frac{\sin x}x\dd x = \frac\pi2$。一致性检验：
$1 - \cos x = 2\sin^2\frac x2$，而代换 $x = 2u$
将 $\int\frac{2\sin^2(x/2)}{x^2}\dd x$ 化为
$\int\frac{\sin^2u}{u^2}\dd u$：两种计算结果一致。

\textbf{11.} $\int_0^\infty\frac{\sin(ax)}x\dd x = \frac\pi2$
对每个 $a > 0$ 成立（作代换 $u = ax$：该积分
在尺度变换下不变）；$\int_\R\eu^{-ax^2}\dd x = \sqrt{\pi/a}$
（作代换 $u = \sqrt a\,x$）。对自变量作尺度变换，
狄利克雷积分保持不变，高斯积分则
除以 $\sqrt a$。

\textbf{12.} 对所有 $t \in [0,1]$ 都成立的控制函数必须
控制 $\sup_{t\in[0,1]}\abs{\eu^{-tx}\frac{\sin
x}x} = \abs{\frac{\sin x}x}$，而它不可积
（见\cref{exo:b3:lebesgue:6}）：控制收敛定理不能越过 $t = 0$。
第 7 题的替代办法 --- 尾部关于参数一致地小，
紧部分由控制收敛定理处理 --- 这是标准的
“一致可积”模式，而
每当条件收敛遇上极限交换时，它便再次出现。

\textbf{13.} 若 $g = \log f$ 为凸函数，则 $f = \exp\circ g$
亦为凸（exp 是凸增函数：$f(\lambda x + (1-\lambda)
y) \leq \eu^{\lambda g(x) + (1-\lambda)g(y)} \leq \lambda f(x)
+ (1-\lambda)f(y)$，最后一步由 exp 在
$g(x), g(y)$ 两点之间的凸性得出）。乘积与仿射代换：
取对数把乘积化为和，凸函数的仿射代换
亦为凸函数。赫尔德不等式（$\frac1p + \frac1q = 1$）：
对 $\int\abs u^p = \int\abs v^q = 1$，由杨不等式得 $\abs{uv} \leq
\frac{\abs u^p}p + \frac{\abs v^q}q$，积分得：$\int\abs{uv}
\leq 1$；一般情形由齐次性推出。再设 $\lambda
\in \intoo01$，取 $p = \frac1\lambda$，对
下列分解式应用上式：
\[
t^{\lambda x + (1-\lambda)y - 1}\eu^{-t}
= \bigl(t^{x-1}\eu^{-t}\bigr)^{\lambda}
\bigl(t^{y-1}\eu^{-t}\bigr)^{1-\lambda} :
\qquad
\Gamma(\lambda x + (1{-}\lambda)y) \leq
\Gamma(x)^\lambda\,\Gamma(y)^{1-\lambda} .
\]

\textbf{14.} 对凸函数 $g$，弦的斜率随
其端点递增（三弦不等式）。比较
区间 $[n-1, n]$、$[n, n+x]$、$[n, n+1]$ 上的弦：
\[
\log(n-1) = \frac{g(n) - g(n-1)}1 \leq
\frac{g(n+x) - g(n)}x \leq \frac{g(n+1) - g(n)}1 = \log n,
\]
乘以 $x > 0$ 即得结论。

\textbf{15.} 设 $f(n) = (n-1)!$（由 $f(1) = 1$ 递推）
且 $f(n + x) = x(x+1)\cdots(x+n-1)\,f(x)$，第 14 题化为
\[
(n-1)^x\,(n-1)! \;\leq\; x(x+1)\cdots(x+n-1)\,f(x)
\;\leq\; n^x\,(n-1)! .
\]
上界可改写为 $f(x) \leq
\frac{n!\,n^x}{x(x+1)\cdots(x+n)}\cdot\frac{x+n}n$，而
取 $n+1$ 时的下界写为 $f(x) \geq
\frac{n!\,n^x}{x(x+1)\cdots(x+n)}$。修正因子
$\frac{x+n}n \to 1$，由夹逼原理可得：
\[
f(x) = \lim_n\frac{n!\,n^x}{x(x+1)\cdots(x+n)}
\qquad (x \in \intoc01),
\]
这一表达式与 $f$ 无关，且在 $\intoc01$ 上唯一，
再由递推关系推至处处。由于 $\Gamma$ 满足
三条假设（第 13 题），故 $f = \Gamma$，高斯公式
成立 --- 对一切 $x > 0$ 皆然，因两边满足同一
递推关系。

\textbf{16.} 在 $0$ 附近，被积函数 $\sim t^{x-1}$，
当且仅当 $x > 0$ 时可积；在 $1$ 附近则对称，$y$ 起同样作用。
在 $\intcc\varepsilon{1-\varepsilon}$ 上分部积分，
令 $\varepsilon \to 0$（当 $x, y >
0$ 时边界项为零）：
\[
B(x{+}1, y) = \Bigl[-t^x\frac{(1-t)^y}y\Bigr]_0^1 +
\frac xy\int_0^1t^{x-1}(1-t)^y\,\dd t
= \frac xy\bigl(B(x, y) - B(x{+}1, y)\bigr),
\]
利用 $(1-t)^y = (1-t)^{y-1}(1 - t)$；解得 $B(x+1, y) =
\frac{x}{x+y}B(x, y)$。又 $B(1, y) =
\int_0^1(1-t)^{y-1}\dd t = \frac1y$。

\textbf{17.} $f(1) = \frac1y\cdot\frac{\Gamma(1+y)}{\Gamma(y)}
= 1$；$f(x+1) = \frac{x}{x+y}B(x,y)\cdot\frac{(x+y)\Gamma(x+y)}
{\Gamma(y)} = x\,f(x)$；且 $f$ 关于 $x$ 对数凸，因为它是
对数凸的 $B(\cdot, y)$（赫尔德不等式作用于
分解式 $t^{(\lambda x_1 + (1-\lambda)x_2)-1}(1-t)^{y-1}
= (\cdots)^\lambda(\cdots)^{1-\lambda}$，同第 13 题）
与 $\Gamma(\cdot + y)$ 之积（仿射平移）。玻尔--莫勒鲁普：$f =
\Gamma$，即 $B(x, y) =
\frac{\Gamma(x)\Gamma(y)}{\Gamma(x+y)}$。

\textbf{18.} 取 $t = \sin^2\theta$，$\dd t =
2\sin\theta\cos\theta\,\dd\theta$ 且 $t^{-1/2}(1-t)^{-1/2} =
\frac1{\sin\theta\cos\theta}$：
\[
B\Bigl(\frac12, \frac12\Bigr) =
\int_0^{\pi/2}2\,\dd\theta = \pi
= \frac{\Gamma(\frac12)^2}{\Gamma(1)}
\quad\Longrightarrow\quad
\Gamma\Bigl(\frac12\Bigr) = \sqrt\pi .
\]
第一部分通过对参数求导，并使两个函数各自趋于极限，
同样求得了这一常数；这里仅凭凸性便将问题固定下来，
直至只剩一个值。
以运动为手段的分析与以形状为手段的分析之对照。

\textbf{19.} $g(1) = \frac{2^0}{\sqrt\pi}\Gamma(\frac12)
\Gamma(1) = 1$。递推关系：
\[
g(x+1) = \frac{2^{x}}{\sqrt\pi}\,
\Gamma\Bigl(\frac{x+1}2\Bigr)\Gamma\Bigl(\frac x2 + 1\Bigr)
= \frac{2^{x}}{\sqrt\pi}\cdot\frac x2\,
\Gamma\Bigl(\frac x2\Bigr)\Gamma\Bigl(\frac{x+1}2\Bigr)
= x\,g(x) .
\]
对数凸性：$\eu^{(x-1)\log2}$（对数仿射）与
对数凸函数 $\Gamma$ 的两个仿射重参数化之积。
玻尔--莫勒鲁普给出 $g = \Gamma$；取 $x = 2z$：
$\Gamma(2z) = \frac{2^{2z-1}}{\sqrt\pi}\Gamma(z)\Gamma(z +
\frac12)$.

\textbf{20.} 由 $\Gamma(\frac12) = \sqrt\pi$ 及递推关系，
recursion, $\Gamma(n + \frac12) = (n - \frac12)(n -
\frac32)\cdots\frac12\,\sqrt\pi =
\frac{(2n-1)(2n-3)\cdots1}{2^n}\sqrt\pi =
\frac{(2n)!}{4^nn!}\sqrt\pi$（用偶数补全
奇数乘积）。渐近估计：由第 14 题，取 $g = \log\Gamma$
及 $x = \frac12$，可得 $\sqrt{n-1} \leq
\frac{\Gamma(n+\frac12)}{\Gamma(n)} \leq \sqrt n$，于是
该比值与 $\sqrt n$ 之商夹在 $\sqrt{1 - \frac1n}$
与 $1$ 之间。

\textbf{21.} 由第 20 题，$(2n)! =
\frac{4^nn!}{\sqrt\pi}\Gamma(n + \tfrac12)$，于是
\[
\binom{2n}n = \frac{(2n)!}{(n!)^2}
= \frac{4^n\,\Gamma(n+\frac12)}{\sqrt\pi\;n!}
= \frac{4^n}{\sqrt\pi}\cdot
\frac{\Gamma(n+\frac12)}{n\,\Gamma(n)}
\sim \frac{4^n}{\sqrt\pi}\cdot\frac{\sqrt n}{n}
= \frac{4^n}{\sqrt{\pi n}} .
\]
数值上，$4^{10}/\sqrt{10\pi} = 1048576/5.6050 \approx
187078.6$，与 $\binom{20}{10} = 184756$ 相比：比值为 $0.9876$
--- 误差为 $O(1/n)$，在 $n = 10$ 时即已可见。

\textbf{22.} 等价关系：$\Gamma(\frac12) = \sqrt\pi
\Leftrightarrow G = \sqrt\pi$（第 3 题的代换 $x = u^2$
）$\Leftrightarrow B(\frac12, \frac12) = \pi$
（欧拉公式）；在 $z = n$ 时的倍元公式\emph{就是}
$\Gamma(n + \frac12)$ 的闭形式，而它\emph{恰是}
在斜率引理的意义下的中心二项式系数估计。参数技巧
（第一、二部分）用于计算运动量的极限，
且在确有真正形变存在时不可或缺（狄利克雷积分
没有凸性证明）；凸性技巧
本身不算出任何东西，却禁止一切
--- 它擅长唯一性问题与函数方程
（高斯、欧拉、勒让德三笔即可），而求导
则会淹没于计算之中。一位完整的
分析学家两者兼用。

\textbf{23.} 取 $t = \sin^2\theta$，$\dd t =
2\sin\theta\cos\theta\,\dd\theta =
2\,t^{1/2}(1-t)^{1/2}\,\dd\theta$，于是
\[
W_p = \int_0^1 t^{p/2}\,
\frac{\dd t}{2\,t^{1/2}(1-t)^{1/2}}
= \frac12\int_0^1 t^{\frac{p+1}2 - 1}(1-t)^{\frac12 - 1}\dd t
= \frac12\,B\Bigl(\frac{p+1}2, \frac12\Bigr).
\]
由贝塔函数的递推关系，取 $x = \frac{n+1}2$，$y = \frac12$
可得
\[
W_{n+2} = \frac12\,
\frac{(n+1)/2}{(n+2)/2}\,B\Bigl(\frac{n+1}2, \frac12\Bigr)
= \frac{n+1}{n+2}\,W_n .
\]
由 $W_0 = \frac\pi2$，$W_1 = 1$ 出发：
\[
W_{2n} = \frac\pi2\prod_{k=1}^n\frac{2k-1}{2k},
\qquad
W_{2n+1} = \prod_{k=1}^n\frac{2k}{2k+1} .
\]
由于 $\sin^{n+1} \leq \sin^n$ 在 $\intcc0{\pi/2}$ 上成立，
数列 $(W_n)$ 单调不增，于是
\[
1 \geq \frac{W_{2n+1}}{W_{2n}} \geq
\frac{W_{2n+1}}{W_{2n-1}} = \frac{2n}{2n+1}
\longrightarrow 1 .
\]
但由闭形式可得
\[
\frac{W_{2n+1}}{W_{2n}} = \frac2\pi
\prod_{k=1}^n\frac{(2k)^2}{(2k-1)(2k+1)}
= \frac2\pi\prod_{k=1}^n\frac{4k^2}{4k^2-1},
\]
令 $n \to \infty$，即得沃利斯乘积。（经
欧拉公式，$W_p = \frac{\sqrt\pi}2\,
\Gamma(\frac{p+1}2)/\Gamma(\frac p2 + 1)$：沃利斯公式
不过是高斯积分换了另一副面孔。）

\textbf{24.} 被积函数关于 $b$ 的导数为
$-2x\,\eu^{-x^2}\sin(2bx)$，被 $2\abs
x\,\eu^{-x^2} \in L^1(\R)$ 控制（关于 $b$ 一致）：$\Phi$ 属于
$\mathcal C^1$，且
\[
\Phi'(b) = -\int_{-\infty}^{+\infty}
2x\,\eu^{-x^2}\sin(2bx)\,\dd x .
\]
分部积分，取 $u = \sin(2bx)$，$\dd v =
-2x\,\eu^{-x^2}\dd x$（故 $v = \eu^{-x^2}$），边界项
为零，于是
\[
\Phi'(b) = -2b\int_{-\infty}^{+\infty}
\eu^{-x^2}\cos(2bx)\,\dd x = -2b\,\Phi(b).
\]
因此 $\bigl(\Phi(b)\,\eu^{b^2}\bigr)' = 0$，且 $\Phi(b) =
\Phi(0)\,\eu^{-b^2} = \sqrt\pi\,\eu^{-b^2}$（由第一部分）。在相差
一个常数因子的意义下，这说明 $\eu^{-x^2}$ 的傅里叶变换
仍是高斯函数——这正是
在\cref{ch:b3:fouriertransform}中反演理论所依托的不动点。

\textbf{25.} 对 $0 < a < b$ 与 $x > 0$，
\[
0 \leq \frac{\eu^{-ax} - \eu^{-bx}}{x}
= \int_a^b \eu^{-sx}\,\dd s \leq (b - a)\,\eu^{-ax},
\]
故积分 $I(a)$ 收敛（勒贝格意义下）；逐点估计
还表明被积函数可连续延拓，延拓值为 $b -
a$，位于 $x = 0$ 处。固定 $b$；在 $\intco{a_0}{+\infty}$ 上，被积函数
关于 $a$ 的导数为 $-\eu^{-ax}$，被
$\eu^{-a_0x} \in L^1(\intoo0{+\infty})$ 控制，故 $I$ 属于 $\mathcal
C^1$，在 $\intoo0b$ 上，且
\[
I'(a) = -\int_0^{+\infty}\eu^{-ax}\dd x = -\frac1a,
\qquad\text{因此}\qquad
I(a) = -\log a + c .
\]
由双侧估计得 $0 \leq I(a) \leq (b-a)/a \to 0$（当
$a \to b^-$ 时），故 $c = \log b$，$I(a) = \log\frac ba$。
可去奇点必须\emph{在 $0$ 处}考虑：每一项
$\eu^{-ax}/x$ 单独在 $0$ 附近的积分是发散的（对数型），
只有一阶消去
$\eu^{-ax} - \eu^{-bx} = O(x)$ 才使差在该处可积；在无穷远处，
每一项本身
都是无害的。
\end{solution}
